% 2010 年上海春季卷第 12 题：程序框图。
% 来源 PDF：春季高考/2010/2010春季上海.pdf，SHA-256 24926e5272a6316cbaa399c0a097b3986ab34563db975016f98bae99f2274c56。
% 原页：第 1 页右栏下部（印刷页 2-37，《数学教学》2010 年第 2 期转载）。
% 重绘依据：开始 → a←√5，b←√7，n←1 → 判断 n(b−a)>1；否则 n←n+1 回到判断；
% 是则 m←[na] → r←(m+1)/n → 输出 r → 结束。循环回线接在判断框上方。
\begin{tikzpicture}[semithick, every node/.style={inner sep=2pt, font=\small}, >=stealth]
  \node[draw, rounded corners] (start) at (0,0) {开始};
  \node[draw, rectangle] (init) at (0,-1.2) {$a\leftarrow\sqrt5,\ b\leftarrow\sqrt7,\ n\leftarrow1$};
  \node[draw, diamond, aspect=2.4] (test) at (0,-3.0) {$n(b-a)>1$};
  \node[draw, rectangle] (inc) at (2.9,-3.0) {$n\leftarrow n+1$};
  \node[draw, rectangle] (m) at (0,-4.4) {$m\leftarrow[na]$};
  \node[draw, rectangle] (r) at (0,-5.4) {$r\leftarrow(m+1)/n$};
  \node[draw, trapezium, trapezium left angle=70, trapezium right angle=110] (out) at (0,-6.5) {输出 $r$};
  \node[draw, rounded corners] (end) at (0,-7.6) {结束};

  \draw[->] (start) -- (init);
  \draw[->] (init) -- (test);
  \draw[->] (test) -- node[above] {否} (inc);
  \draw[->] (inc.north) -- (2.9,-2.05) -- (0,-2.05);
  \draw[->] (0,-2.05) -- (test.north);
  \draw[->] (test) -- node[right] {是} (m);
  \draw[->] (m) -- (r);
  \draw[->] (r) -- (out);
  \draw[->] (out) -- (end);
\end{tikzpicture}
